\chapter{Grafovska predstavitev podatkov sledilnika pogleda}
\label{sec:grafovska-predstavitev}

To poglavje je eno osrednjih metodoloških poglavij naloge.
Namen je natančno opisati, kako iz časovnega okna meritev sledilnika pogleda nastane graf, ki ga uporabimo kot vhod v GNN.

\TODO{
\begin{itemize}
    \item To poglavje naj bo edino mesto za natančno definicijo konkretnega vhodnega grafa te naloge.
    \item Ne ponavljaj splošnih definicij grafa, heterogenosti ali uteženosti iz poglavja~\ref{sec:ozadje}; uporabi sklic.
    \item Ne opisuj arhitekturnih odločitev modela, kot so relacijske transformacije, MLP za uteži, normalizacija naučenih uteži, \texttt{GCNConv} ali READOUT; te prestavi v poglavje~\ref{sec:modeli}.
\end{itemize}
}

\TODO{
\begin{itemize}
    \item Dodati opombo, da se grafovska konstrukcija v eksperimentih prilagodi izbranemu naboru signalov.
    \item Za \emph{gaze} opisati graf iz koordinat pogleda $(x,y)$ in relacij, ki jih je smiselno zgraditi iz pogleda; za \emph{gaze+pupil} dodati zenične značilke vozlišč; za \emph{all} ohraniti trenutno polno konstrukcijo z razdaljo, fiksacijami, fiksacijskimi povezavami in značilkami povezav.
    \item \emph{Pupil-only} ne vključevati v glavno grafovsko konstrukcijo, dokler ni jasen smiseln GNN protokol; zaenkrat ga označiti kot odprto možnost ali ga izpustiti.
\end{itemize}
}

\TODO{
\begin{itemize}
    \item Glavno besedilo: dodati eno osrednjo metodološko sliko grafovske predstavitve. Slika naj na majhnem, berljivem izseku okna pokaže vozlišča po koordinatah pogleda in relacije časovnih, prostorskih ter fiksacijskih povezav z različnimi barvami ali paneli.
    \item Glavno besedilo: slika naj razloži idejo konstrukcije grafa, ne vseh implementacijskih podrobnosti. Ne dodajaj ločene slike za vsako relacijo v glavnem besedilu, ker bi to prekinilo tok poglavja.
    \item Dodatek~\ref{app:grafi}: prestaviti podrobne vizualizacije posameznih relacij, primer 10-sekundnega grafa, berljiv 2- ali 5-sekundni izsek, prekrivanja relacij in različice z velikostjo vozlišč po zenici.
\end{itemize}
}

%%%%%%%%%%%%%
% Uredniški komentar za prepis poglavja:
% - Poglavje 5 naj odgovori na vprašanje: kaj je vhodni graf oziroma katera grafovska predstavitev nastane iz enega 10-sekundnega okna.
% - Poglavje 6 naj odgovori na vprašanje: kako GNN ta graf obdela. Tja sodijo posredovanje sporočil, relacijske transformacije, GCNConv, MLP za uteži povezav, normalizacija uteži, READOUT, klasifikacijska glava in izguba.
% - V tem poglavju pusti samo kratke napovedi tipa "te značilke povezav uporabi model, opisan v poglavju~\ref{sec:modeli}", brez formul za naučene uteži.
% - Prejšnji TODO glede pupil-only popravi: samo zenici je ena od štirih ekvivalentnih GNN konstrukcij. Graf je smiseln že zaradi časovnih povezav; nima pa prostorskih in fiksacijskih relacij, če množica signalov ne vsebuje koordinat pogleda oziroma fiksacijske informacije.
%%%%%%%%%%%%%

\section{Učni vzorec kot graf}
\label{subsec:ucni-vzorec}

\subsection{Definicija grafa}

Vsak učni vzorec je časovno okno meritev sledilnika pogleda.
Okno pretvorimo v graf
\begin{equation}
    G = (V, E_t^+, E_t^-, E_s, E_f),
\end{equation}
kjer je $V$ množica vozlišč, $E_t^+$ množica časovnih povezav naprej, $E_t^-$ množica časovnih povezav nazaj, $E_s$ množica prostorskih povezav in $E_f$ množica fiksacijskih povezav.

%%%%%%%%%%%%%
% Kaj tukaj spremeniti:
% - Definicijo grafa napiši tako, da najprej poda splošno obliko za polno konstrukcijo, nato pa pove, da je dejanska množica povezav odvisna od izbrane množice signalov.
% - Za vse štiri množice signalov mora nastati graf:
%   - samo pogled: vozlišča imajo čas in koordinate pogleda; relacije so časovne in prostorske;
%   - samo zenici: vozlišča imajo čas in velikosti zenic; relacije so časovne;
%   - pogled in zenici: vozlišča imajo čas, koordinate pogleda in velikosti zenic; relacije so časovne in prostorske;
%   - vsi signali: vozlišča imajo še oddaljenost od sledilnika pogleda in fiksacijsko informacijo; relacije so časovne, prostorske in fiksacijske.
% - Razmisli o zapisu z množico relacij $\mathcal{R}_{S}$, kjer $S$ označuje izbrano množico signalov, namesto da polna oblika $G=(V,E_t^+,E_t^-,E_s,E_f)$ izgleda kot edina možna oblika.
% - Ne uvajaj tukaj modelnih imen; to je še vedno definicija vhodnega učnega vzorca.
%%%%%%%%%%%%%

Vozlišča predstavljajo zaporedne meritve sledilnika pogleda znotraj časovnega okna.
Povezave med vozlišči opisujejo različne relacije med meritvami: časovno zaporednost, prostorsko bližino pogleda in pripadnost isti fiksaciji.
En graf predstavlja en učni vzorec.

\subsection{Oznaka grafa}

\TODO{
\begin{itemize}
    \item Uskladi z novo formulacijo naloge: glavni cilj je binarna klasifikacija valence.
    \item Tu zapiši samo, da je oznaka na ravni grafa oziroma časovnega okna; podrobna preslikava oznak naj ostane v poglavju~\ref{sec:podatki} oziroma \ref{sec:eksperimenti}.
\end{itemize}
}

Vsak graf označimo z oznako pripadajočega časovnega okna.
V nalogi to oznako uporabimo pri ločenih klasifikacijskih nalogah valence in vzburjenosti.
\END_NOTE{Na koncu popravi, če klasifikacijski nalogi nista ločeni, ampak združeni v eno nalogo. Predvidevam, da bosta ločeni, ker mislim, da so tako tudi naredili v MAHNOB članku.}

%%%%%%%%%%%%%
% Kaj tukaj spremeniti:
% - Zamenjaj stavek o ločenih nalogah valence in vzburjenosti.
% - Povej samo, da graf dobi oznako binarne valence na ravni časovnega okna.
% - Dodaj, da je oznaka podedovana od kombinacije udeleženec-posnetek, zato imajo vsa okna istega ogleda isti ciljni razred.
% - Ne razlagaj tukaj preslikave iz čustvenih oznak v valenco; sklic na poglavje~\ref{sec:podatki} oziroma dodatek je dovolj.
%%%%%%%%%%%%%


\section{Vozlišča}
\label{subsec:vozlisca}

\subsection{Definicija vozlišča}

Vsako vozlišče predstavlja eno meritev sledilnika pogleda znotraj časovnega okna:
\begin{equation}
    v_i \leftrightarrow \mathbf{x}_i.
\end{equation}
Če okno vsebuje $n$ veljavnih meritev, potem graf vsebuje $n$ vozlišč.
Ker so meritve vzorčene s frekvenco 60~Hz, je pričakovano število meritev v polnem oknu
\begin{equation}
    n_{\mathrm{exp}} = 10~s \cdot 60~Hz = 600.
\end{equation}
% oceni, ali je bolje ta podatek podati tukaj ali v poglavju 4, kjer je to tudi že omenjeno. ne želimo se ponavljati, zato to omenimo samo na enem mestu, na drugem mestu pa popolnoma zbrišemo.

Grafovska predstavitev naravno podpira spremenljivo število vozlišč.

\subsection{Značilke vozlišč}
\label{subsec:znacilke-vozlisca}

\TODO{
\begin{itemize}
    \item Tu naj ostane samo opis vhodnih značilk vozlišč kot dela grafovske predstavitve.
    \item Preslikavo značilk vozlišč v latentni prostor opiši v poglavju~6.4.
    \item Ne opisuj LayerNorm, dropout, MLP blokov ali drugih arhitekturnih komponent v tem poglavju.
\end{itemize}
}

Vektor značilk vozlišča je
\begin{equation}
    \mathbf{x}_i =
    \left[
        x_i,
        y_i,
        p_{l,i},
        p_{r,i},
        t_i,
        d_i,
        f_{dur,i}
    \right],
\end{equation}
kjer sta $x_i$ in $y_i$ koordinati pogleda, $p_{l,i}$ in $p_{r,i}$ velikosti leve in desne zenice, $t_i$ je normalizirani čas znotraj okna, $d_i$ razdalja med sledilnikom in udeleženčevimi očmi, $f_{dur,i}$ pa trajanje fiksacije, ki ji pripada $i$-ta meritev.


\subsubsection{Časovna značilka}

Časovna značilka $t_i$ modelu poda položaj meritve znotraj okna.

\subsubsection{Koordinate pogleda}

Koordinati $x_i$ in $y_i$ opisujeta položaj pogleda in sta uporabljeni tudi pri konstrukciji prostorskih povezav.

\subsubsection{Značilke zenice}

Značilki $p_{l,i}$ in $p_{r,i}$ opisujeta velikost leve in desne zenice.

\subsubsection{Razdalja do sledilnika in fiksacije}

Značilki $d_i$ in $f_{dur,i}$ opisujeta razdaljo do sledilnika ter trajanje fiksacije pripadajoče meritve.
Manjkajoče vrednosti trajanja fiksacije in meritve, ki ne pripadajo fiksaciji, so kodirane z vrednostjo $f_{dur,i}=0$.

\section{Heterogena grafovska predstavitev}
\label{subsec:heterogeni-graf}

\TODO{
\begin{itemize}
    \item Splošna definicija heterogenega grafa naj ostane v poglavju~\ref{sec:ozadje}.
    \item Tukaj definiraj samo konkretno množico relacij $\mathcal{R}$ in razloži, kaj posamezna relacija pomeni v podatkih sledilnika pogleda.
    \item Razlago, da model za vsako relacijo uporablja ločeno transformacijo oziroma ločeno posredovanje sporočil, prestavi v poglavje~\ref{subsec:relacijsko-posredovanje}.
\end{itemize}
}

\subsection{Tipi relacij}

Ker različne povezave opisujejo različne odnose med meritvami, graf obravnavamo kot heterogen graf z več tipi povezav.
V trenutni predstavitvi uporabljamo množico relacij
\begin{equation}
    \mathcal{R} =
    \{\texttt{temporal\_forward}, \texttt{temporal\_backward}, \texttt{spatial}, \texttt{fixation}\}.
\end{equation}
Časovne povezave naprej, časovne povezave nazaj, prostorske povezave in fiksacijske povezave so obravnavane kot ločene relacije.

%%%%%%%%%%%%%
% Kaj tukaj spremeniti:
% - Množica relacij naj ne bo opisana kot vedno enaka za vse eksperimente.
% - Zapiši polno množico relacij za vsi signali, nato pa povej, da se pri drugih množicah uporabi podmnožica:
%   - samo pogled: časovne relacije naprej/nazaj in prostorska relacija;
%   - samo zenici: samo časovne relacije naprej/nazaj;
%   - pogled in zenici: časovne relacije naprej/nazaj in prostorska relacija;
%   - vsi signali: časovne relacije naprej/nazaj, prostorska relacija in fiksacijska relacija.
% - Tu lahko ostane konceptualna utemeljitev ločevanja relacij, ampak brez stavkov o tem, da model uporablja ločeno transformacijo. To prestavi v poglavje~\ref{subsec:relacijsko-posredovanje}.
%%%%%%%%%%%%%

\subsection{Zakaj ločimo relacije}

Ločevanje relacij omogoča, da model za vsak tip relacije uporablja ločeno transformacijo oziroma ločeno funkcijo posredovanja sporočil.
S tem lahko časovno zaporednost, prostorsko bližino in pripadnost isti fiksaciji modeliramo z različnimi parametri, namesto da bi vse povezave obravnavali kot enakovredne. 
S tem tudi dovolimo modelu, da se za različne relacije nauči dajati rezlično težo istim vozliščem. Morda je neko vozlišče zelo pomembno s pogleda časovne zaporednosti, ampak povsem nepomembno z vidika umeščenosti v prostor.

%%%%%%%%%%%%%
% Kaj tukaj spremeniti:
% - Odstrani modelno razlago "ločena transformacija", "funkcija posredovanja sporočil" in "nauči dajati različno težo"; to sodi v poglavje 6.
% - Tukaj zapiši samo podatkovno motivacijo: časovna povezava pomeni zaporednost meritev, prostorska povezava pomeni bližino položajev pogleda, fiksacijska povezava pomeni pripadnost isti zaznani fiksaciji.
% - Popravi tipkarsko napako "rezlično" in formulacijo "s pogleda"; še bolje pa celoten odstavek zamenjaj s krajšo vsebinsko motivacijo.
%%%%%%%%%%%%%

\TODO{
\begin{itemize}
    \item Ta pododsek naj ostane konceptualna motivacija za heterogen graf.
    \item Podrobnosti, kako model dejansko uporabi ločene relacijske transformacije, prestavi v poglavje~6.4.
    \item Popravi slog in tipkarske napake; izogni se preveč modelno specifični razlagi v poglavju 5.
\end{itemize}
}

\section{Povezave}

\subsection{Časovne povezave}
\label{subsec:casovne-povezave}

Časovne povezave povezujejo meritve zaporedne v času.
Naj bodo vozlišča v posameznem oknu urejena po času, tako da je $v_i$ $i$-ta meritev v časovnem zaporedju.
Parameter $k_t \in \mathbb{N}$ določa število časovnih sosedov na vsako stran.
V privzeti nastavitvi uporabljamo $k_t = 2$. \END_NOTE{preveri končni $k_t$}

%%%%%%%%%%%%%
% Kaj tukaj spremeniti:
% - Časovne povezave veljajo za vse štiri množice signalov, tudi za samo zenici.
% - Preveri končno konfiguracijo: v MEMORY piše izbrana nastavitev `kt=1`, starejše besedilo pa omenja $k_t=2$.
% - Ko bo nastavitev potrjena, naj bodo $k_t$ v besedilu, formulah, slikah in tabeli statistike grafov skladni.
%%%%%%%%%%%%%

\subsubsection{Povezave naprej in nazaj}

Ločimo usmerjene časovne povezave naprej in nazaj:
\begin{align}
    E_t^+ &= \{(v_i, v_j) : 1 \leq j - i \leq k_t\}, \\
    E_t^- &= \{(v_i, v_j) : 1 \leq i - j \leq k_t\}.
\end{align}

Povezave v $E_t^+$ potekajo od prejšnjih proti kasnejšim meritvam, povezave v $E_t^-$ pa od kasnejših proti prejšnjim meritvam.
Ločitev smeri je smiselna, ker ima vrstni red očesnega gibanja pomen.
Primer časovnih povezav na majhnem izseku okna je prikazan na sliki~\ref{fig:casovne-povezave-primer}.

\begin{figure}[H]
    \centering
    \includegraphics[width=\textwidth]{gnn_graph_far_02_P10_90_avi_t106_86-107_95_kt1_ks1_kf3_f3_n9/casovne.png}
    \caption{Primer časovnih povezav v izseku grafa. Povezave sledijo časovni urejenosti meritev in ločijo smer naprej ter nazaj.}
    \label{fig:casovne-povezave-primer}
\end{figure}

%%%%%%%%%%%%%
% Kaj tukaj spremeniti:
% - To ločeno sliko najverjetneje prestavi v dodatek~\ref{app:grafi}.
% - V glavnem besedilu naj ostane ena osrednja slika, ki skupaj pokaže časovne, prostorske in fiksacijske povezave.
% - Če slika ostane v glavnem besedilu, jo uporabi samo, če res pomaga razložiti konstrukcijo in ne podvaja osrednje slike.
%%%%%%%%%%%%%

\subsubsection{Parameter $k_t$}

Parameter $k_t$ določa velikost lokalnega časovnega konteksta in s tem število povezav, ki jih dodamo okrog posamezne meritve. Smiselno lahko zahtevamo $k_t < n$, kjer je $n$ število vozlišč v oknu. Pričakovano število usmerjenih povezav v grafu naprej je $|E_t^+| = k_t \cdot n - \sum_{i=1}^{k_t} (i) = k_t \cdot n - \frac{k_t(k_t+1)}{2} = k_t \cdot (n - \frac{k_t+1}{2})$. Za obravnavani primer \END_NOTE{$k_t=2$ in $n=600$ je to $|E_t^+| = 2 \cdot (600 - \frac{2+1}{2}) = 1197$ }usmerjenih povezav naprej, enako število pa je tudi usmerjenih povezav nazaj. Torej skupno imamo približno $2 k_t \cdot n$ časovnih povezav.

\END_NOTE*{
    \begin{itemize}
        \item V glavno metodološko sliko vključi časovne povezave, v dodatek~\ref{app:grafi} pa prestavi ločeni vizualizaciji \texttt{temporal\_forward} in \texttt{temporal\_backward}.
        \item Razmisli o smiselnosti $k_t$ glede na število layerjev GNN .
    \end{itemize}
}


\subsection{Prostorske povezave}
\label{subsec:prostorske-povezave}

Prostorske povezave so neusmerjene povezave, ki povezujejo meritve s podobnim položajem pogleda znotraj istega časovnega okna.
Za vsako vozlišče poiščemo $k_s$ najbližjih sosedov v prostoru koordinat pogleda $(x,y)$.

%%%%%%%%%%%%%
% Kaj tukaj spremeniti:
% - Jasno zapiši, da prostorske povezave obstajajo samo pri množicah signalov, ki vsebujejo koordinate pogleda: samo pogled, pogled in zenici, vsi signali.
% - Pri množici samo zenici prostorskih povezav ne gradimo, ker v tej predstavitvi ni koordinat pogleda.
% - Preveri končno konfiguracijo: v MEMORY piše `ks=1`, starejše besedilo in tabela pa omenjata $k_s=2$.
%%%%%%%%%%%%%

\subsubsection{Najbližji sosedi v prostoru pogleda}

Prostorsko razdaljo med meritvama $i$ in $j$ definiramo evklidsko:
\begin{equation}
    d_{ij}^{(s)} =
    \sqrt{(x_i - x_j)^2 + (y_i - y_j)^2}.
\end{equation}

Za vsako vozlišče $v_i$ poiščemo $k_s$ najbližjih sosedov v prostoru koordinat pogleda $(x,y)$ z uporabo k-d drevesa. Če je $v_j$ med temi sosedi, dodamo povezavo med vozliščema $v_i$ in $v_j$; ker prostorske povezave obravnavamo kot neusmerjene, povezavo predstavimo z obema usmerjenima povezavama $(v_i,v_j)$ in $(v_j,v_i)$. Na koncu odstranimo podvojene povezave znotraj relacije $E_s$.

\subsubsection{Simetrizacija povezav}

Ker povezave obravnavamo kot neusmerjene, vsaki povezavi $(v_i,v_j)$ dodamo tudi povezavo $(v_j,v_i)$.
Primer prostorskih povezav je prikazan na sliki~\ref{fig:prostorske-povezave-primer}.

\begin{figure}[H]
    \centering
    \includegraphics[width=\textwidth]{gnn_graph_far_02_P10_90_avi_t106_86-107_95_kt1_ks1_kf3_f3_n9/prostorske.png}
    \caption{Primer prostorskih povezav v izseku grafa. Povezave nastanejo med meritvami s podobnim položajem pogleda, zato niso nujno omejene na časovno zaporedne meritve.}
    \label{fig:prostorske-povezave-primer}
\end{figure}

%%%%%%%%%%%%%
% Kaj tukaj spremeniti:
% - To ločeno sliko najverjetneje prestavi v dodatek~\ref{app:grafi}.
% - V glavnem besedilu naj prostorske povezave pokrije osrednja metodološka slika.
%%%%%%%%%%%%%

\END_NOTE*{
    \begin{itemize}
        \item V glavno metodološko sliko vključi prostorske povezave, v dodatek~\ref{app:grafi} pa prestavi ločeno vizualizacijo \texttt{spatial} povezav.
        \item $k_s=2$ v trenutnih poskusih. Preveri končni $k_s$.
    \end{itemize}
}


\subsection{Fiksacijske povezave}
\label{subsec:fiksacijske-povezave}

Fiksacijske povezave povezujejo meritve, ki pripadajo isti zaznani fiksaciji.
Pripadnost fiksaciji določimo z identifikatorjem fiksacije, podanim v \MAHNOB podatkih.

%%%%%%%%%%%%%
% Kaj tukaj spremeniti:
% - Jasno zapiši, da fiksacijske povezave obstajajo samo pri množici vsi signali, ker samo ta vključuje fiksacijsko informacijo.
% - Pri samo pogled in pogled in zenici se lahko koordinate uporabijo za prostorske povezave, ne pa za fiksacijske povezave, če fiksacijska informacija ni del te množice signalov.
% - Pri samo zenici ostanejo samo časovne povezave.
%%%%%%%%%%%%%

\subsubsection{Redčenje povezav znotraj fiksacije}

Znotraj posamezne fiksacije ne povežemo vseh vozlišč med seboj, saj bi to povzročilo preveliko število fiksacijskih povezav.
Namesto tega uporabimo redkejšo konstrukcijo, ki ohrani informacijo o pripadnosti isti fiksaciji, vendar zmanjša prostorsko in časovno zahtevnost konstrukcije grafa ter učenja na grafu.
To konstrukcijo imenujemo razširjene povezave znotraj fiksacije (angl. \emph{dilated intra-fixation edges}).
Fiksacijo obravnavamo kot zvezno zaporedje vozlišč z istim veljavnim identifikatorjem \texttt{fixation-index}.
Če ima taka fiksacija $F$ vozlišč, jim znotraj fiksacije priredimo lokalne indekse $0,\ldots,F-1$.

Gostoto povezav uravnava parameter $k_f$, ki določa število fiksacijskih povezav, ki jih generira vsako vozlišče.
V trenutni privzeti konfiguraciji uporabljamo $k_f=3$.
Za posamezno fiksacijo najprej izračunamo korak razširitve
\begin{equation}
    s = \max(1, \lfloor F/k_f + 0{,}5 \rfloor).
\end{equation}
Množico kandidatnih odmikov nato definiramo kot
\begin{equation}
    D = \{(1 + q s) \bmod F \mid q = 0,\ldots,k_f-1\}.
\end{equation}
Morebitni odmik $0$ in podvojene odmike odstranimo.
Za vsak odmik $d \in D$ dodamo neusmerjene povezave
\begin{equation}
    \{v_i, v_{(i+d)\bmod F}\},
    \qquad i = 0,\ldots,F-1.
\end{equation}
Pri zapisu grafa za posredovanje sporočil vsako neusmerjeno povezavo predstavimo z dvema usmerjenima povezavama.
Podvojene povezave znotraj relacije $E_f$ odstranimo.

Na primer, pri fiksaciji dolžine $F=30$ in parametru $k_f=3$ dobimo $s=10$.
Vozlišče z lokalnim indeksom $0$ torej povežemo z vozlišči $1, 11$ in $21$.
Odmik $1$ ohrani časovno lokalno komunikacijo znotraj fiksacije, večji odmiki pa omogočijo širjenje informacij tudi med časovno bolj oddaljenimi meritvami iste fiksacije.
Primer fiksacijskih povezav je prikazan na sliki~\ref{fig:fiksacijske-povezave-primer}.

\begin{figure}[H]
    \centering
    \includegraphics[width=\textwidth]{gnn_graph_far_02_P10_90_avi_t106_86-107_95_kt1_ks1_kf3_f3_n9/fiksacijske.png}
    \caption{Primer fiksacijskih povezav v izseku grafa. Povezave združujejo meritve, ki pripadajo isti zaznani fiksaciji, pri čemer razširjena konstrukcija ohrani redkejšo povezovalno strukturo.}
    \label{fig:fiksacijske-povezave-primer}
\end{figure}

%%%%%%%%%%%%%
% Kaj tukaj spremeniti:
% - To ločeno sliko najverjetneje prestavi v dodatek~\ref{app:grafi}.
% - Če želiš v glavnem besedilu razložiti razširjene povezave znotraj fiksacije, naj bo to s kratkim primerom v besedilu ali kot panel v osrednji sliki.
%%%%%%%%%%%%%

Razširjene povezave imajo večji doseg kot zgolj zaporedne fiksacijske povezave, hkrati pa se izognejo kombinatorični eksploziji, ki bi nastala pri polnem kliku.
Pri polnem kliku bi namreč število povezav naraščalo kvadratno z dolžino fiksacije.
Pri analiziranih 10-sekundnih oknih dobimo približno $1.956$ fiksacijskih povezav na graf pri $k_f=2$, $2.908$ pri $k_f=3$, $4.163$ pri $k_f=5$ in $7.273$ pri $k_f=10$.
Pri privzeti izbiri $k_f=3$ so fiksacijske povezave predstavljale približno $45{,}7\,\%$ vseh povezav v grafu (relacije \texttt{temporal}, \texttt{spatial} in \texttt{fixation}).

Izbira $k_f=3$ je zato kompromis med dosegom in gostoto grafa.
Ker vsaka plast GNN omogoča izmenjavo informacij med neposredno povezanimi vozlišči, se doseg znotraj fiksacije povečuje tudi z večanjem števila plasti.
Pri $k_f=3$ vozlišče že v eni plasti prejme informacije iz več delov iste fiksacije, večplastni model pa lahko te informacije dodatno propagira po fiksacijski relaciji.
Večje vrednosti $k_f$ povečajo gostoto grafa, računsko zahtevnost in možnosti za prekomerno glajenje, manjše vrednosti pa graf dodatno razredčijo.

%%%%%%%%%%%%%
% Kaj tukaj spremeniti:
% - Odstrani ali prestavi stavke o številu plasti GNN, propagaciji po plasteh in prekomernem glajenju; to je razlaga obdelave grafa v modelu.
% - Tukaj pusti samo konstrukcijsko utemeljitev: polni klik znotraj fiksacije je pregost, zato uporabimo redkejšo povezovalno shemo, ki ohrani pripadnost isti fiksaciji.
% - Če želiš ohraniti kompromis $k_f$, ga zapiši z vidika gostote grafa in števila povezav, ne z vidika delovanja večplastnega GNN.
%%%%%%%%%%%%%
\END_NOTE*{
    \begin{itemize}
        \item V glavno metodološko sliko vključi fiksacijske povezave, v dodatek~\ref{app:grafi} pa prestavi ločeno vizualizacijo \texttt{fixation} povezav in primer razširjenih povezav znotraj ene fiksacije.
        \item $k_f=3$ v trenutnih poskusih. Preveri končni $k_f$ glede na eksperimente (consider kf=2).
    \end{itemize}
}

Isti izsek z vsemi opisanimi tipi povezav je prikazan na sliki~\ref{fig:konstrukcija-grafa-primer}.
Prikaz povzema konstrukcijo heterogenega grafa in zato služi kot povezava med formalno definicijo grafa ter posameznimi konstrukcijskimi pravili.

\begin{figure}[H]
    \centering
    \includegraphics[width=\textwidth]{gnn_graph_far_02_P10_90_avi_t106_86-107_95_kt1_ks1_kf3_f3_n9/vse_povezave.png}
    \caption{Majhen primer konstrukcije grafa iz podatkov sledilnika pogleda. Vozlišča predstavljajo meritve v izseku časovnega okna, povezave pa skupaj prikazujejo časovne, prostorske in fiksacijske relacije, opisane v tem poglavju.}
    \label{fig:konstrukcija-grafa-primer}
\end{figure}

%%%%%%%%%%%%%
% Kaj tukaj spremeniti:
% - To naj postane glavna oziroma edina slika konstrukcije grafa v poglavju 5.
% - Napis naj omeni, da slika prikazuje polno konstrukcijo za množico vsi signali; druge množice signalov uporabljajo ustrezne podmnožice značilk in relacij.
% - Preveri, ali pot datoteke ustreza končni lokaciji v `slike/konstrukcija_grafa` in ali je slika v dogovorjeni paleti B.
%%%%%%%%%%%%%

\section{Uteži povezav}
\label{subsec:utezi-povezav}

\TODO{
\begin{itemize}
    \item Ta odsek preimenuj v \emph{Značilke povezav}, če bodo naučene uteži in njihova normalizacija prestavljene v poglavje~\ref{sec:modeli}.
    \item Razdeli vsebino tega odseka med poglavji 5 in 6.4.
    \item V poglavju 5 naj ostanejo značilke povezav: katere količine opišejo relacijo med dvema meritvama.
    \item V poglavje~6.4 prestavi ali tam povzemalno opiši arhitekturni del: MLP za izračun uteži, predznačene uteži, normalizacijo in uporabo uteži pri posredovanju sporočil.
    \item V poglavju 5 ne poudarjaj več, da model iz značilk izračuna uteži; to naj bo predvsem del opisa predlaganega modela.
\end{itemize}
}

Uteži povezav omogočajo, da model razlikuje med povezavami, ki povezujejo v času in prostoru različno oddaljene meritve.

\subsection{Značilke povezav}

%%%%%%%%%%%%%
% Kaj tukaj spremeniti:
% - Razdelek preimenuj v "Značilke povezav".
% - Prvi stavek naj ne govori, da uteži omogočajo modelu razlikovanje; to je arhitekturna razlaga iz poglavja 6.
% - Tukaj povej samo, da nekaterim povezavam poleg tipa relacije priredimo še opisne značilke, ki povzamejo časovno, prostorsko in signalno razliko med povezanima meritvama.
% - Pojasni, katere značilke obstajajo pri katerih množicah signalov:
%   - časovne značilke so na voljo pri vseh množicah;
%   - $\Delta x$, $\Delta y$ in $d_{ij}^{screen}$ so na voljo samo, kadar imamo koordinate pogleda;
%   - $\Delta d^{eye}_{ij}$ je na voljo samo pri vsi signali;
%   - smer $\rho_{ij}$ je smiselna samo pri časovnih povezavah.
% - Ne piši formule $a_{ij}=f_{\theta}(\mathbf{r}_{ij})$ v tem poglavju; prestavi jo v poglavje~\ref{subsec:relacijsko-posredovanje}.
%%%%%%%%%%%%%

Za vsako povezavo sestavimo vektor značilk, iz katerega model izračuna surovo skalarno oceno povezave:
\begin{equation}
    a_{ij}=f_{\theta}(\mathbf{r}_{ij}),
\end{equation}
kjer je $\mathbf{r}_{ij}$ vektor značilk povezave $e_{ij}$, $f_{\theta}$ pa večslojni perceptron.
Za časovne, prostorske in fiksacijske povezave so uporabljene značilke $[t_i, t_j, \Delta t, \Delta x, \Delta y, d_{ij}^{screen}, \Delta d^{eye}_{ij}]$, pri časovnih povezavah pa dodatno še značilka smeri $\rho_{ij}$.
Pri časovni povezavi naprej je $\rho_{ij}=1$, pri časovni povezavi nazaj pa $\rho_{ij}=-1$.
Časovne povezave imajo torej 8 značilk, prostorske in fiksacijske povezave pa 7 značilk, iz katerih se izračunajo uteži povezav.
  
\begin{table}[H]
    \centering
    \caption{Značilke povezav za izračun naučenih uteži povezav.}
    \label{tab:znacilke-povezav}
    \begin{tabularx}{\textwidth}{@{}>{\raggedright\arraybackslash}p{0.22\textwidth} >{\raggedright\arraybackslash}X@{}}
        \toprule
        Značilka & Opis \\
        \midrule
        $t_i$ & Normaliziran čas izvornega vozlišča znotraj časovnega okna. \\
        $t_j$ & Normaliziran čas ciljnega vozlišča znotraj časovnega okna. \\
        $\Delta t$ & Časovna razlika med ciljnim in izvornim vozliščem, $\Delta t = t_j - t_i$. \\
        $\Delta x$ & Razlika v vodoravni koordinati pogleda, $\Delta x = x_j - x_i$. \\
        $\Delta y$ & Razlika v navpični koordinati pogleda, $\Delta y = y_j - y_i$. \\
        $d_{ij}^{screen}$ & Evklidska razdalja med položajema pogleda na zaslonu. \\
        $\Delta d^{eye}_{ij}$ & Sprememba razdalje med sledilnikom in udeleženčevimi očmi. \\
        $\rho_{ij} = \pm 1$ & Smer časovne povezave; uporablja se samo pri časovnih povezavah. \\
        \bottomrule
    \end{tabularx}
\end{table}

%%%%%%%%%%%%%
% Kaj tukaj spremeniti:
% - Napis tabele spremeni tako, da ne govori "za izračun naučenih uteži povezav", ampak samo "Značilke povezav".
% - V tabeli označi, pri katerih relacijah oziroma množicah signalov je posamezna značilka uporabljena, ali pa to pojasni v besedilu pred tabelo.
% - Če tabela postane preširoka, jo razdeli na stolpce "Značilka", "Opis", "Uporaba".
%%%%%%%%%%%%%


\subsection{Normalizacija uteži}

\TODO{
\begin{itemize}
    \item Razmisli o prestavitvi tega pododseka v poglavje~6.4, ker je normalizacija uteži del arhitekture modela in ne del same definicije vhodnega grafa.
    \item Če formule ostanejo tukaj, jih v 6.4 ne ponavljaj; tam dodaj samo konceptualni povzetek in sklic na ta odsek.
    \item Če formule prestaviš v 6.4, naj poglavje 5 vsebuje samo tabelo značilk povezav.
\end{itemize}
}

Surovo skalarno oceno $a_{ij}$ najprej preslikamo v omejeno območje na interval $(-1, 1)$:
\begin{equation}
    \tilde{w}_{ij} = \tanh(a_{ij}).
\end{equation}

Nato uteži normaliziramo po vhodnih povezavah ciljnega vozlišča. 
To pomeni, da za vozlišče $v_j$ upoštevamo povezave, ki imajo $v_j$ za ciljno vozlišče.
Pri neusmerjenih povezavah, tretiramo vse povezave za vhodne in izhodne.
Normalizacijo izvedemo ločeno za vsako relacijo $r$:
\begin{equation}
    w_{ij}^{(r)} =
    \frac{\tilde{w}_{ij}^{(r)}}
    {\sum_{i' \in \mathcal{N}^{-}_r(j)}
    |\tilde{w}_{i'j}^{(r)}| + \epsilon},
\end{equation}
kjer je $\mathcal{N}^{-}_r(j)$ množica izvornih vozlišč, ki so z vozliščem $v_j$ povezana z vhodno povezavo tipa $r$.
Majhna konstanta $\epsilon$ prepreči deljenje z nič.
Takšna normalizacija stabilizira velikost agregiranih sporočil tudi pri vozliščih z večjim številom vhodnih povezav.

Za izračun surovih ocen uporabljamo tri ločene večslojne perceptrone: enega za časovne povezave, enega za prostorske povezave in enega za fiksacijske povezave.
Časovni relaciji naprej in nazaj si delita isti perceptron, ločeni pa sta z dodatno značilko smeri $\rho_{ij}$.
Skalarne uteži povezav se uporabljajo pri konvoluciji \texttt{GCNConv}, ki je privzeta izbira modela.
Pri poskusih z \texttt{GATConv} se te skalarne uteži ne uporabijo, saj \texttt{GATConv} pomembnost povezav modelira z notranjim mehanizmom pozornosti.

%%%%%%%%%%%%%
% Kaj tukaj spremeniti:
% - Celoten pododsek "Normalizacija uteži" prestavi v poglavje~\ref{subsec:relacijsko-posredovanje}.
% - V poglavju 5 lahko ostane največ en stavek, da so značilke povezav vhod v uteževanje povezav pri modelu HeteroGCN-MLP-w, ki je opisan v poglavju~\ref{sec:modeli}.
% - Omembe \texttt{GCNConv}, \texttt{GATConv}, MLP-jev za uteži in normalizacije po vhodnih povezavah ne sodijo v definicijo grafovske predstavitve.
%%%%%%%%%%%%%

\section{Statistika grafov}
\label{subsec:statistika-grafov}

\TODO{
\begin{itemize}
    \item Ta razdelek lahko ostane v poglavju~\ref{sec:grafovska-predstavitev}, ker opisuje lastnosti konstruiranih grafov.
    \item V njem ne razlagaj več splošne računske zahtevnosti GNN modela; za to uporabi poglavje~\ref{sec:modeli}, poglavje~\ref{sec:eksperimenti} ali dodatek s hiperparametri.
    \item Preveri, da so parametri v napisu tabele skladni s končno izbrano konfiguracijo in da niso v konfliktu z opisom v poglavju~\ref{sec:eksperimenti}.
\end{itemize}
}

\subsection{Število vozlišč in povezav}

\begin{table}[H]
    \centering
    \caption{Povzetek statistike grafov za privzeto 10-sekundno okno pri $k_t=2$, $k_s=2$ in $k_f=3$.}
    \label{tab:graph-statistics}
    \begin{tabular}{lr}
        \toprule
        Količina & Vrednost \\
        \midrule
        Mediana števila vozlišč & 560 \\
        Povprečno število vozlišč & 505,9 \\
        Mediana števila časovnih povezav naprej & 1117 \\
        Mediana števila časovnih povezav nazaj & 1117 \\
        Mediana števila prostorskih povezav & 1546 \\
        Mediana števila fiksacijskih povezav & 3204 \\
        Povprečna gostota grafa & TODO \\
        \bottomrule
    \end{tabular}
\end{table}

%%%%%%%%%%%%%
% Kaj tukaj spremeniti:
% - Statistiko preračunaj za končno konfiguracijo in za vsako množico signalov posebej, ker samo zenici, samo pogled, pogled in zenici ter vsi signali nimajo istih relacij.
% - Tabela naj verjetno vsebuje stolpce po množicah signalov in vrstice: število vozlišč, časovne povezave naprej, časovne povezave nazaj, prostorske povezave, fiksacijske povezave, skupno število povezav.
% - Za relacije, ki v določeni množici signalov ne obstajajo, uporabi pomišljaj ali nič, vendar pojasni izbiro.
% - Preveri skladnost z izbrano konfiguracijo iz MEMORY: `kt=1`, `ks=1`, `kf=3`, razen če koda pokaže drugače.
%%%%%%%%%%%%%

\END_NOTE{Verjetno prestavi to poglavje nižje med modele.}
 
\subsection{Računska zahtevnost}

\TODO{
\begin{itemize}
    \item Ta odsek verjetno ne sodi v poglavje 5, ker govori o zahtevnosti modela, ne o konstrukciji vhodnega grafa.
    \item Statistiko števila vozlišč in povezav lahko pustiš v poglavju 5.
    \item Oceno števila parametrov, časa učenja, pomnilniške zahtevnosti in primerjavo z drugimi modeli prestavi v poglavje 6, poglavje 7 ali dodatek s hiperparametri.
    \item Če v poglavju 5 ostane kaj o zahtevnosti, naj bo omejeno na zahtevnost konstrukcije grafa in gostoto povezav.
\end{itemize}
}

% TODO: Izpolni spodnje alineje in preveri podatke
\TODO{Stvari spodaj}
\begin{itemize}
    \item Oceni velikost GNN modela po neki metriki -- število učljivih parametrov.
    \item Primerjaj velikost med modeli:
    \begin{itemize}
        \item GNN približno 560K parametrov, MLP baseline približno 5K, GazeMAE + MLP približno 2.100K parametrov, od tega samo 80K učljivih (MLP), ker je GazeMAE hrbtenica zamrznjena. (Te podatke je potrebno preveriti.)
        \item LightGBM lahko merimo v n\_estimators in max\_depth, ampak ni zelo primerljivo.
        \item SVM je odvisen od števila podpornih vektorjev. SVM in LGBM bomo zato primerjali samo časovno.
    \end{itemize}
    \item Pri eni plasti GNN je strošek posredovanja sporočil približno linearen v številu povezav in v širini skritih predstavitev. Za $L$ plasti je zato groba odvisnost $O(L \cdot |E| \cdot H)$, ob tem pa model vsebuje še linearne transformacije in MLP-je za združevanje relacij. Pri predlaganem modelu je $L=3$, $H=128$, tipičen graf pa ima nekaj sto vozlišč in nekaj tisoč povezav.
    \item Primerjaj prostorsko zahtevnost med modeli.
    \item Primerjaj časovno zahtevnost med modeli -- konkreten povprečen čas učenja za vsak model v tej nalogi.
    \item Verjetno prestavi to poglavje nižje med modele.
\end{itemize}

%%%%%%%%%%%%%
% Kaj tukaj spremeniti:
% - Ta odsek odstrani iz poglavja 5.
% - Primerjavo parametrov, časa učenja in časa inference prestavi v poglavje 7 ali v rezultate/praktično zahtevnost, ker primerja modele in ne grafovske predstavitve.
% - Če želiš v poglavju 5 ohraniti kaj o zahtevnosti, naj bo omejeno na konstrukcijo grafa: kako število povezav raste s $k_t$, $k_s$, $k_f$ in zakaj ne uporabimo polnega klika znotraj fiksacije.
%%%%%%%%%%%%%
